%-------------------------------------------------------------------------------
%  CV — Daniel González Pascual
%  Estilo: NIT Warangal Resume (Overleaf gtsrbjvffcjn)
%  Compilar con pdfLaTeX en Overleaf (o texlive local).
%-------------------------------------------------------------------------------
\documentclass[letterpaper,11pt]{article}

\usepackage{latexsym}
\usepackage[empty]{fullpage}
\usepackage{titlesec}
\usepackage{marvosym}
\usepackage[usenames,dvipsnames]{color}
\usepackage{verbatim}
\usepackage{enumitem}
\usepackage[hidelinks]{hyperref}
\usepackage{fancyhdr}
\usepackage[spanish,es-noshorthands]{babel}
\usepackage[utf8]{inputenc}
\usepackage[T1]{fontenc}
\usepackage{tabularx}
\usepackage{fontawesome5}
\input{glyphtounicode}

\pagestyle{fancy}
\fancyhf{}
\fancyfoot{}
\renewcommand{\headrulewidth}{0pt}
\renewcommand{\footrulewidth}{0pt}

% Márgenes ajustados al estilo de la plantilla
\addtolength{\oddsidemargin}{-0.5in}
\addtolength{\evensidemargin}{-0.5in}
\addtolength{\textwidth}{1in}
\addtolength{\topmargin}{-0.5in}
\addtolength{\textheight}{1.0in}

\urlstyle{same}

\raggedbottom
\raggedright
\setlength{\tabcolsep}{0in}

% Secciones: versalitas + regla (firma visual de la plantilla NIT Warangal)
\titleformat{\section}{
  \vspace{-4pt}\scshape\raggedright\large
}{}{0em}{}[\color{black}\titlerule\vspace{-5pt}]

\pdfgentounicode=1

%---- Comandos auxiliares -----------------------------------------------------
\newcommand{\resumeItem}[1]{
  \item\small{#1\vspace{-2pt}}
}
\newcommand{\resumeSubheading}[4]{
  \vspace{-2pt}\item
    \begin{tabular*}{0.97\textwidth}[t]{l@{\extracolsep{\fill}}r}
      \textbf{#1} & #2 \\
      \textit{\small#3} & \textit{\small #4} \\
    \end{tabular*}\vspace{-7pt}
}
\newcommand{\resumeProjectHeading}[2]{
  \item
    \begin{tabular*}{0.97\textwidth}{l@{\extracolsep{\fill}}r}
      \small#1 & #2 \\
    \end{tabular*}\vspace{-7pt}
}
\newcommand{\resumeSubItem}[1]{\small#1\vspace{-4pt}}
\newcommand{\resumeSubList}{\begin{itemize}[leftmargin=0.15in, label={}]}
\newcommand{\resumeEndSubList}{\end{itemize}}
\newcommand{\resumeItemListStart}{\begin{itemize}}
\newcommand{\resumeItemListEnd}{\end{itemize}\vspace{-5pt}}

%-------------------------------------------------------------------------------
\begin{document}

%---------- CABECERA ----------
\begin{center}
  {\Huge \scshape Daniel González Pascual} \\ \vspace{4pt}
  \small Desarrollador Backend Python $\cdot$ Semi-Senior \\ \vspace{6pt}
  \small
  \faPhone\ 625 382 700 \;
  \faEnvelope\ \href{mailto:danigpascual@protonmail.com}{danigpascual@protonmail.com} \;
  \faMapMarker\ Málaga, España \;
  \faGithub\ \href{https://github.com/danigpas}{github.com/danigpas} \;
  \faLinkedin\ \href{https://www.linkedin.com/in/daniel-gonz%C3%A1lez-pascual-dev/}{linkedin.com/in/daniel-gonzález-pascual-dev}
\end{center}

%---------- SOBRE MÍ ----------
\section{Sobre mí}
\small{
Desarrollador Backend Python con más de 3 años construyendo APIs escalables e integraciones en tiempo
real (FastAPI, Celery, Redis, RabbitMQ). Me gusta entender el sistema completo: del diseño de la API
al pipeline de datos y al despliegue — mantengo mi propio homelab con Proxmox y Docker, donde
publico proyectos como Portfolio OS, un escritorio funcional en el navegador. Actualmente trabajo
en proyectos del sector energético europeo con Python, PySpark y AWS. Meticuloso, proactivo y
obsesionado con la trazabilidad, el rendimiento y que nada se caiga.
}

%---------- EXPERIENCIA ----------
\section{Experiencia}
\resumeSubList

  \resumeSubheading
    {Inforyde}{Mar 2026 -- Actualidad}
    {Desarrollador Python Semi-Senior}{Madrid, España}
    \resumeItemListStart
      \resumeItem{Backend en \textbf{Python} para una comercializadora suiza de electricidad y gas: desarrollo de servicios y procesos de datos sobre distintos proyectos del sector energético.}
      \resumeItem{Procesamiento de grandes volúmenes de datos con \textbf{PySpark} sobre \textbf{AWS}.}
      \resumeItem{Integración y despliegue continuo con \textbf{Docker} y \textbf{Jenkins}, gestión ágil con Jira y trabajo en equipos multidisciplinares.}
    \resumeItemListEnd

  \resumeSubheading
    {DisOfic SLU}{Ago 2023 -- Mar 2026}
    {Desarrollador Backend}{Málaga, España}
    \resumeItemListStart
      \resumeItem{Diseño e implementación de una API en tiempo real con \textbf{FastAPI} que sincroniza \textbf{+50 sitios WordPress} con el CRM \textbf{Odoo}: stock y precios de +12.000 productos visibles en las páginas de producto y gestión automatizada de pedidos.}
      \resumeItem{Automatización de procesos de datos con workers \textbf{Celery / Redis / RabbitMQ}: sincronización masiva de clientes y artículos reducida de días a 2 horas.}
      \resumeItem{Reducción del \textbf{50\%} en incidencias de pedidos gracias a la sincronización en tiempo real.}
      \resumeItem{Administración y optimización de consultas sobre \textbf{MySQL, PostgreSQL y Oracle}.}
    \resumeItemListEnd

\resumeEndSubList

%---------- PROYECTOS ----------
\section{Proyectos}
\resumeSubList

  \resumeProjectHeading
    {\textbf{Portfolio OS} — \emph{Next.js 15, TypeScript, FastAPI, Docker}}{2026}
    \resumeItemListStart
      \resumeItem{Escritorio Omarchy funcional en el navegador con ventanas reales, terminal y backend FastAPI con fallback offline-first. Publicado en \href{https://danigpascual.dev}{danigpascual.dev}.}
    \resumeItemListEnd

  \resumeProjectHeading
    {\textbf{API de integración WordPress \ensuremath{\leftrightarrow} Odoo} — \emph{FastAPI, Celery, Redis, RabbitMQ}}{2023 -- 2026}
    \resumeItemListStart
      \resumeItem{Sincronización en tiempo real de +50 sitios con el CRM: pedidos, stock y sincronización masiva de clientes y artículos.}
    \resumeItemListEnd

  \resumeProjectHeading
    {\textbf{El Nieto de Pascual} — \emph{Blog}}{2024 -- Actualidad}
    \resumeItemListStart
      \resumeItem{Blog personal sobre desarrollo backend, Python y aprendizajes construyendo proyectos.}
    \resumeItemListEnd

\resumeEndSubList

%---------- EDUCACIÓN ----------
\section{Educación}
\resumeSubList

  \resumeSubheading
    {Junta de Andalucía}{Sept 2026 -- Jun 2028}
    {Grado Superior en Administración de Sistemas Informáticos en Red (ASIR) -- en curso}{Málaga, España}
  \resumeSubheading
    {Universidad de Málaga}{2017 -- 2021}
    {Grado en Ingeniería Informática}{Málaga, España}
  \resumeSubheading
    {Cesur Málaga}{Ene -- Ago 2023}
    {Certificado de Profesionalidad IFCT0609: Programación de Sistemas Informáticos}{Málaga, España}

\resumeEndSubList

%---------- HABILIDADES ----------
\section{Habilidades}
\resumeSubList

  \resumeProjectHeading{\textbf{Lenguajes}}{Python (FastAPI, Flask, Django, Reflex), C\#, SQL, PySpark}
  \resumeProjectHeading{\textbf{Bases de datos}}{PostgreSQL, MySQL, Oracle}
  \resumeProjectHeading{\textbf{Cloud \& CI/CD}}{AWS, Jenkins, Docker, Jira}
  \resumeProjectHeading{\textbf{Backend \& mensajería}}{Celery, Redis, RabbitMQ, Pydantic, REST}
  \resumeProjectHeading{\textbf{ERP \& CMS}}{Odoo, WordPress}
  \resumeProjectHeading{\textbf{Automatización}}{Apache Airflow, n8n}
  \resumeProjectHeading{\textbf{Herramientas}}{Git, GitHub, GitLab, Linux, Proxmox}
  \resumeProjectHeading{\textbf{Web}}{TypeScript, Next.js, React, Tailwind CSS}

\resumeEndSubList

\end{document}
